% META: {"id": "TSPE-PROBA-CR-010", "chapitre": "TSPE-PROBABILITES", "type_objet": "cours", "capacites_codes": ["C1"], "sources_inspiration": [], "mode_creation": "ex_nihilo", "status": "approved"}

\section{Succession d'épreuves, arbres pondérés}

% ============================================================
% STRATE 1 — Rappels et méthode
% ============================================================

\propriete[Règles de l'arbre pondéré]{
\begin{itemize}
  \item Sur chaque branche issue d'un même noeud, la somme des probabilités vaut $1$.
  \item La probabilité d'un chemin est le \textbf{produit} des probabilités des branches qui le composent.
  \item La probabilité d'un événement est la \textbf{somme} des probabilités des chemins qui le realisent.
\end{itemize}
}

\definition[1]{
Deux epreuves sont \textbf{indépendantes} si le résultat de l'une n'influence pas les probabilités de l'autre. Deux événements $A$ et $B$ sont indépendants si $P(A \cap B) = P(A) \times P(B)$.
}

\propriete[Formule des probabilités totales]{
Si $A_1, \ldots, A_n$ forment une partition de l'univers (événements disjoints dont l'union est l'univers, tous de probabilité non nulle), alors pour tout événement $B$ :
\[
  P(B) = \sum_{i=1}^n P(A_i) \times P_{A_i}(B).
\]
}

\exemple{
Une urne contient $3$ boules rouges et $2$ boules vertes. On tire deux boules successivement, sans remise. Quelle est la probabilité de tirer deux boules de couleurs différentes ?

Arbre : $P(R_1) = \dfrac{3}{5}$, puis $P_{R_1}(V_2) = \dfrac{2}{4}$ ; $P(V_1)=\dfrac{2}{5}$, puis $P_{V_1}(R_2)=\dfrac{3}{4}$.
\[
  P(\text{couleurs differentes}) = \frac{3}{5}\times\frac{2}{4} + \frac{2}{5}\times\frac{3}{4} = \frac{6}{20}+\frac{6}{20} = \frac{12}{20} = \frac{3}{5}.
\]

% BEGIN-VERIFY
% from sympy import Rational
% p = Rational(3,5)*Rational(2,4) + Rational(2,5)*Rational(3,4)
% assert p == Rational(3,5)
% END-VERIFY
}

% ============================================================
% STRATE 2 — Erreurs frequentes
% ============================================================

\erreurFrequente{
\textbf{Confondre tirage avec remise et sans remise.} Sans remise, les probabilités conditionnelles changent a chaque étape (le dénominateur diminue). Avec remise, les probabilités restent identiques a chaque étape (epreuves indépendantes et identiques).
}
